% ===========================================================================
\chapter{Rules as joins}\label{ch:joins}

A rule's head holds whenever its body does, so applying a rule means
finding every way to make its body true. Take

\begin{shell}
in_common(X, Y) :- follows(X, Z), follows(Y, Z), X != Y.
\end{shell}

The body has three conditions. Match the first, \code{follows(X, Z)}, against
every \code{follows} fact; each match binds \code{X} and \code{Z}. For each
of those, match the second, \code{follows(Y, Z)}, where \code{Z} now has a
value, so only the other followers of that same person match, binding
\code{Y}. For each of those, check the test \code{X != Y}. Every set of
bindings that survives all three gives one fact, \code{in_common(X, Y)}.

Combining conditions this way, each one narrowing the possibilities the
earlier ones left, is what a relational database calls a \emph{join}, and
it is the heart of the engine.

\section{The join}

The join goes through the body one literal at a time, carrying the
bindings so far. It is recursive: to satisfy literals \code{i} onwards,
satisfy literal \code{i} in every possible way, and for each way, satisfy
literals \code{i + 1} onwards. When it runs off the end of the body, the
bindings satisfy everything, and they are a solution.

It belongs to a class, \code{Evaluator}, because the next chapters will need
two kinds of evaluator that join in slightly different ways. The class also
counts \code{attempts}, the facts the join has tried to match: a simple
measure of how much work an evaluation does, which we will use to compare
evaluators.

\begin{newfile}{evaluate.nex}
-- evaluate.nex: deriving every fact the rules imply.

intern Database

-- A way of bringing a database to a fixpoint of some rules: a state in
-- which applying any rule derives nothing new.
deferred class Evaluator
feature
  -- How many facts the joins have looked at: a measure of the work done.
  attempts: Integer

  -- Every set of bindings that satisfies the body of `r` in `db`.
  solutions(r: Rule, db: Database): Array[Bindings]
  do
    result := []
    join(r.body, 0, create Bindings.empty, db, result)
  end

  -- Extend `b` to satisfy body.get(i), body.get(i + 1) and so on, adding
  -- every complete extension to `found`.
  join(body: Array[Literal], i: Integer, b: Bindings, db: Database, found: Array[Bindings])
  require
    in_body: 0 <= i and i <= body.length
  do
    if i = body.length then
      found.add(b)
    else
      match body.get(i) of
        Positive(atom) then
          across source(i, atom, b, db) as f do
            attempts := attempts + 1
            let next := match_atom(atom, f, b)
            if next /= nil then
              join(body, i + 1, next, db, found)
            end
          end
        Compare(op, left, right) then
          if holds(op, substitute_term(left, b), substitute_term(right, b)) then
            join(body, i + 1, b, db, found)
          end
        Negated(atom) then
          if not db.contains(substitute(atom, b)) then
            join(body, i + 1, b, db, found)
          end
      end
    end
  end

  -- Where literal `i` of a body finds its facts: here, the database.
  source(i: Integer, atom: Atom, b: Bindings, db: Database): Array[Atom]
  do
    result := db.candidates(atom, b)
  end
end
\end{newfile}

Each kind of literal is handled differently:
\begin{itemize}
  \item An atom is matched against every candidate fact, and the join
    continues with each match.
  \item A test, \code{Compare}, is checked with the bindings applied; if it
    holds, the join continues with the bindings unchanged.
  \item A negated atom, \code{not banned(U)}, holds when the fact, with the
    bindings applied, is \emph{not} in the database. Whether that is a sound
    thing to ask is the subject of Chapter~\ref{ch:negation}.
\end{itemize}

Where an atom's candidate facts come from is decided by \code{source}, which
for now asks the database. Giving that decision a routine of its own is
what will let the faster evaluator of Chapter~\ref{ch:seminaive} change it.

The order of the body matters for tests. A test can only be checked once its
variables have values, so \code{X != Y} must come after the atoms that bind
\code{X} and \code{Y}. Chapter~\ref{ch:safety} turns that into a rule the
engine checks.

\section{Trying it out}

\code{Evaluator} is deferred, so the try program makes a small subclass of
its own just to call \code{solutions}. It finds every pair of people who
follow someone in common:

\begin{tryit}{try_joins.nex}
intern Parser
intern io/Path
intern Evaluate

class Joiner inherit Evaluator
create
  make
  do
  end
end

match parse_program(create Path.make(program.working_directory + "/social.dl").read_text) of
  Ok as ok then
    do
      let db := create Database.make
      across ok.value.facts as f do
        db.add(f)
      end
      let in_common := ok.value.rules.get(3)
      console.print_line(in_common.text)
      let j := create Joiner.make
      across j.solutions(in_common, db) as b do
        console.print_line("  " + b.text + "   gives " + substitute(in_common.head, b).text)
      end
      console.print_line("facts looked at: " + j.attempts.to_string)
    end
  Err as e then console.print_line(e.error)
end
\end{tryit}
\tryoutput{try_joins}

Bo and Fay both follow Cy, so the pair appears, twice, once each way round,
because \code{in_common} is symmetric. Nobody is paired with themselves,
thanks to \code{X != Y}: without it, everyone who follows anyone would have
something in common with themselves.
